\section{数据预处理}

\subsection{数据结构整理}

题目附件同时包含地块、作物、历史种植、统计数据和结果模板。若直接把表格送入优化模型，合并单元格、说明行、价格区间文本和缺失组合都会导致变量索引不统一。因此，本文先把数据整理成可被模型直接调用的四类对象：地块表、作物表、2023 年种植锚点、作物参数表。

附件1的耕地表共有54块地。平旱地、梯田、山坡地、水浇地、普通大棚和智慧大棚分别保留其原始面积，露天耕地总面积核对为1201亩，与题面一致。附件1的作物表前41行为有效作物信息，说明性文字只用于建立种植规则，不参与数值计算。根据题目规则，本文构造可行种植集合 $\Omega$：只有当地块类型、作物类型和季次同时满足种植制度时，三元组 $(i,j,s)$ 才进入模型。该步骤可以提前剔除不可能种植的组合，减少后续优化规模。

附件2的2023年种植情况表中，部分地块名因合并单元格显示为空。本文采用前向填充恢复地块名，并将“单季”统一记为第一季。这样处理后，可得到 $u^{2023}_{i,j,s}$，用于刻画2024年与2023年的连作关系，也用于问题三中前茬豆类的判定。附件2的统计表中，销售单价以区间形式给出，本文取区间中点作为基准价格；亩产量和种植成本按“作物--地块类型--季次”查表。

\subsection{参数构造}

优化模型最核心的输入包括亩产、成本、价格和预期销售量。亩产、成本、价格分别记为 $Y_{i,j,s}$、$C_{i,j,s}$ 和 $P_{i,j,s}$，由作物、地块类型和季次共同决定。由于优化变量按地块编号展开，本文将地块类型 $k(i)$ 代入统计表，得到每一个可行三元组 $(i,j,s)\in\Omega$ 的参数值。

预期销售量 $D_{j,s}$ 用2023年同作物同季次的全村实际产量估计：
\begin{equation}
  D_{j,s}=\sum_{i\in I} a^{2023}_{i,j,s}Y_{i,j,s},
\end{equation}
其中 $a^{2023}_{i,j,s}$ 表示2023年地块 $i$ 在季次 $s$ 种植作物 $j$ 的面积。这样做的原因是题目没有额外给出市场需求量，而2023年的实际种植和销售统计最能反映当地基准经营规模。

附件2缺少智慧大棚第一季蔬菜的亩产和成本记录，但题目允许智慧大棚第一季种植蔬菜。若直接删除这些组合，会削弱智慧大棚在模型中的作用；若自行设定新参数，又缺乏数据依据。本文采用普通大棚第一季同作物数据进行平补，使智慧大棚第一季保留可行性，同时不额外夸大其产量优势。价格差异仍按统计表中已有口径体现。

\subsection{数据侧写}

图\ref{fig:data-land}展示了不同地块类型的面积分布。露天耕地贡献了主要面积，其中梯田和平旱地占比较高；普通大棚和智慧大棚面积虽小，但因可种植蔬菜和食用菌，在高价值作物安排中具有较强约束价值。

\begin{figure}[H]
  \centering
  \begin{subfigure}{0.48\textwidth}
    \centering
    \includegraphics[width=\textwidth]{data_land_area.png}
    \caption{六类地块面积}
  \end{subfigure}
  \begin{subfigure}{0.48\textwidth}
    \centering
    \includegraphics[width=\textwidth]{data_land_area_pie.png}
    \caption{六类地块面积占比}
  \end{subfigure}
  \caption{地块面积结构}
  \label{fig:data-land}
\end{figure}

图\ref{fig:data-price}给出了41种作物基准销售价格的区间均值。食用菌类作物价格明显高于普通粮食和多数蔬菜，说明高价格作物可能带来更高边际收益；但这些作物通常受到大棚季次、销售量和轮作安排约束，不能仅按价格高低直接扩大种植。

\begin{figure}[H]
  \centering
  \includegraphics[width=0.92\textwidth]{data_price_range.png}
  \caption{41种作物销售价格区间均值}
  \label{fig:data-price}
\end{figure}
